\documentclass[11pt]{article}
\usepackage[a4paper,margin=24mm]{geometry}
\usepackage[T1]{fontenc}
\usepackage{lmodern}
\usepackage{amsmath,amssymb,mathtools,amsthm,bm}
\usepackage{microtype}
\usepackage[hidelinks]{hyperref}
\newcommand{\R}{\mathbb{R}}
\newcommand{\bx}{\bm{x}}
\newcommand{\by}{\bm{y}}
\newcommand{\bz}{\bm{z}}
\newcommand{\bk}{\bm{k}}
\newcommand{\bzeta}{\boldsymbol{\zeta}}
\newcommand{\bkappa}{\boldsymbol{\kappa}}
\newcommand{\dd}{\,\mathrm{d}}
\newtheorem{proposition}{Proposition}
\allowdisplaybreaks[1]
\setlength{\emergencystretch}{2em}

\title{Background velocity estimation by focusing optimization\\
for a smooth layered medium}
\author{}
\date{}

\begin{document}
\maketitle

\section{Geometry, model, and scope}

We consider the three-dimensional geometry of the original imaging
function: the transverse variables belong to $\R^2$, and the longitudinal
coordinate is positive below the array. The true and trial velocities are
smooth, uniformly positive, bounded profiles depending only on depth:
\begin{equation}
c_0=c_0(x_\parallel),\qquad c=c(x_\parallel).
\label{eq:profiles}
\end{equation}
The profiles are extended smoothly across the array plane when defining
the full-space Green's functions. There is no reflecting boundary at
$x_\parallel=0$. We use the scalar constant-density Helmholtz model
\begin{equation}
\left[\Delta_{\by}+\frac{\omega^2}{c(y_\parallel)^2}\right]
\widehat G(\omega,(\bx_\perp,0),\by)
=-\delta\big(\by-(\bx_\perp,0)\big),
\label{eq:helmholtz}
\end{equation}
with outgoing conditions. The function $\widehat G_0$ is defined in the
same way with $c_0$. We retain the direct outgoing arrival in the
high-frequency, slowly varying regime. Smoothness alone does not make
this approximation exact: wavelength-scale variations and significant
reflected arrivals are excluded from the leading model below.

The single-scattering data model is
\begin{align}
\widehat u(\omega,\bx_n,\bx_{n'})
={}&\frac{1}{B}\widehat f\!\left(\frac{\omega-\omega_0}{B}\right)
\int_{\R^2\times\R^+}\rho(\by)
\widehat G_0(\omega,(\bx_n,0),\by)
\widehat G_0(\omega,\by,(\bx_{n'},0))\dd\by.
\label{eq:data}
\end{align}
The reflectivity is supported below the array and away from the source
plane. We do not assume that it varies slowly. Reciprocity is used for
the second Green's function. Our Fourier convention is
\begin{equation}
\widehat f(\omega)=\int_\R f(t)e^{i\omega t}\dd t,
\qquad
f(t)=\frac{1}{2\pi}\int_\R\widehat f(\omega)e^{-i\omega t}\dd\omega.
\label{eq:fourier}
\end{equation}
The spectrum is concentrated near the positive carrier $\omega_0$, with
$B\ll\omega_0$, and $f$ is localized around zero.

As in the original notation, let
\begin{equation}
\bz=(\bz_\perp,z_\parallel),\qquad
\bz_\pm=(\bz_\perp\pm\bzeta_\perp/2,z_\parallel).
\label{eq:zpm}
\end{equation}
The two-point imaging function is
\begin{align}
\mathcal I(\bz,\bzeta_\perp;c)
={}&\int_\R\dd\omega
\iint_{\R^2\times\R^2}\dd\bx_n\dd\bx_{n'}
e^{-(|\bx_n|^2+|\bx_{n'}|^2)/a^2}
\nonumber\\
&\times
\overline{\widehat G(\omega,(\bx_n,0),\bz_+)}
\widehat u(\omega,\bx_n,\bx_{n'})
\overline{\widehat G(\omega,(\bx_{n'},0),\bz_-)}.
\label{eq:image}
\end{align}
All subsequent factorizations are leading-order statements in a reduced
model. They are not exact identities for \eqref{eq:helmholtz}.
The assumption \eqref{eq:profiles} is essential: an arbitrary smooth
map depending on the transverse coordinates requires a different ray
geometry and transverse phase, and is not covered by this derivation.

\section{Travel time and amplitude for the smooth profile}

\subsection{Step 1: derive the paraxial travel time}

The eikonal equation is $|\nabla T|^2=c^{-2}$. For a direct ray from the
origin to $(\bx_\perp,x_\parallel)$, stratification conserves transverse
slowness. Its magnitude is the existing variable $p$, with
$p\sup_{0\leq s\leq x_\parallel}c(s)<1$. The ray representation is
\begin{align}
T(\bx_\perp,x_\parallel)
&=p|\bx_\perp|
+\int_0^{x_\parallel}\sqrt{c(s)^{-2}-p^2}\dd s,
\label{eq:ray-time}\\
|\bx_\perp|
&=p\int_0^{x_\parallel}
\frac{c(s)}{\sqrt{1-p^2c(s)^2}}\dd s.
\label{eq:ray-arrival}
\end{align}
The second identity also follows by stationarity of the first expression
with respect to $p$.

Retain the definitions
\begin{equation}
T_c(x_\parallel)=\int_0^{x_\parallel}\frac{\dd s}{c(s)},
\qquad
U_c(x_\parallel)=\int_0^{x_\parallel}c(s)\dd s.
\label{eq:TU}
\end{equation}
The definitions of $T_0,U_0$ use $c_0$ instead of $c$.
The typical speed $\overline c$ is assumed comparable to both profiles
on the propagation region. Expansion of \eqref{eq:ray-arrival} gives
\begin{equation}
|\bx_\perp|
=pU_c(x_\parallel)
+\frac{p^3}{2}\int_0^{x_\parallel}c(s)^3\dd s
+O\big(p^5\overline c^{\,5}x_\parallel\big).
\label{eq:p-expansion}
\end{equation}
Inverting this relation, first at linear order and then at cubic order,
yields
\begin{align}
p={}&\frac{|\bx_\perp|}{U_c(x_\parallel)}
-\frac{|\bx_\perp|^3}{2U_c(x_\parallel)^4}
\int_0^{x_\parallel}c(s)^3\dd s
+O\!\left(\frac{|\bx_\perp|^5}{\overline c\,x_\parallel^5}\right).
\label{eq:p-inversion}
\end{align}
Indeed, the cubic correction to $pU_c$ must cancel the cubic term in
\eqref{eq:p-expansion}. Expanding the square root in \eqref{eq:ray-time}
before substituting \eqref{eq:p-inversion} gives
\begin{align}
T(\bx_\perp,x_\parallel)
={}&T_c(x_\parallel)+p|\bx_\perp|
-\frac{p^2}{2}U_c(x_\parallel)
-\frac{p^4}{8}\int_0^{x_\parallel}c(s)^3\dd s
+O\big(p^6\overline c^{\,5}x_\parallel\big)\\
={}&T_c(x_\parallel)
+\frac{|\bx_\perp|^2}{2U_c(x_\parallel)}
-\frac{|\bx_\perp|^4}{8U_c(x_\parallel)^4}
\int_0^{x_\parallel}c(s)^3\dd s
+O\!\left(\frac{|\bx_\perp|^6}{\overline c\,x_\parallel^5}\right).
\label{eq:travel-fourth}
\end{align}
Consequently, the quadratic paraxial travel time is
\begin{equation}
T(\bx_\perp,x_\parallel)
\simeq T_c(x_\parallel)
+\frac{|\bx_\perp|^2}{2U_c(x_\parallel)}.
\label{eq:travel-paraxial}
\end{equation}
Besides small ray angles, using this approximation inside an oscillatory
Green's function requires the discarded phase to be small:
\begin{equation}
\frac{\omega_0|\bx_\perp|^4}{U_c(x_\parallel)^4}
\int_0^{x_\parallel}c(s)^3\dd s\ll1.
\label{eq:phase-remainder}
\end{equation}
This requirement applies to the true and trial propagation paths over the
effective aperture and offset range.

\subsection{Step 2: derive the Green's function amplitude}

Writing the leading Green's function as $A e^{i\omega T}$ and
substituting it into \eqref{eq:helmholtz} away from the source gives
\begin{equation}
\Delta A+i\omega\big(2\nabla A\cdot\nabla T+A\Delta T\big)
+\omega^2 A\big(c^{-2}-|\nabla T|^2\big)=0.
\end{equation}
The terms of orders $\omega^2$ and $\omega$ yield, respectively,
\begin{equation}
|\nabla T|^2=c^{-2},\qquad
2\nabla A\cdot\nabla T+A\Delta T=0,
\end{equation}
or equivalently $\nabla\cdot(A^2\nabla T)=0$.

To compute the amplitude explicitly, evaluate the paraxial transport
equation on the central ray. From \eqref{eq:travel-paraxial},
\begin{equation}
\partial_{x_\parallel}T(\mathbf0,x_\parallel)
=\frac{1}{c(x_\parallel)},\qquad
\Delta T(\mathbf0,x_\parallel)
=-\frac{c'(x_\parallel)}{c(x_\parallel)^2}
+\frac{2}{U_c(x_\parallel)}.
\end{equation}
The factor $2$ comes from the two transverse coordinates. Taking the
leading amplitude to depend on depth gives
\begin{equation}
\frac{2A'}{c}+A\left(-\frac{c'}{c^2}+\frac{2}{U_c}\right)=0,
\qquad
\frac{2A'}{A}=\frac{c'}{c}-\frac{2c}{U_c}.
\end{equation}
Since $U_c'=c$, integration gives
\begin{equation}
\frac{A^2}{c(x_\parallel)}U_c(x_\parallel)^2
=\text{constant along the ray}.
\end{equation}
This also expresses conservation of flux in a paraxial ray tube whose
cross-sectional area is proportional to $U_c(x_\parallel)^2$.
Close to the source, $U_c(x_\parallel)\sim c(0)x_\parallel$ and the
leading paraxial homogeneous amplitude is $1/(4\pi x_\parallel)$.
The constant is thus $c(0)/(16\pi^2)$, so that
\begin{equation}
A\simeq\frac{\sqrt{c(0)c(x_\parallel)}}{4\pi U_c(x_\parallel)}.
\label{eq:amplitude}
\end{equation}
This also specifies the model dependence of the amplitude: the same
prefactor should not be assumed for a variable-density operator.

Transverse translation invariance now gives
\begin{align}
\widehat G_0(\omega,(\bx_n,0),\by)
&\simeq\frac{\sqrt{c_0(0)c_0(y_\parallel)}}{4\pi U_0(y_\parallel)}
\exp\!\left\{i\omega\left[T_0(y_\parallel)
+\frac{|\bx_n-\by_\perp|^2}{2U_0(y_\parallel)}\right]\right\},
\label{eq:G0}\\
\widehat G(\omega,(\bx_n,0),\bz_\pm)
&\simeq\frac{\sqrt{c(0)c(z_\parallel)}}{4\pi U_c(z_\parallel)}
\exp\!\left\{i\omega\left[T_c(z_\parallel)
+\frac{|\bx_n-\bz_\perp\mp\bzeta_\perp/2|^2}{2U_c(z_\parallel)}\right]\right\}.
\label{eq:Gc}
\end{align}

\section{Point-source data: derive the offset profile}

\subsection{Step 3: choose the depth coordinate and matched true depth}

Fix $\bz_\perp$ and $u_\parallel>0$. Positivity of the profiles makes
$T_c,U_c,T_0,U_0$ strictly increasing. Define, as in the original text,
\begin{align}
z_{\parallel,c}&=U_c^{-1}(u_\parallel),
&y_{\parallel,c}&=T_0^{-1}\big(T_c(z_{\parallel,c})\big),
\label{eq:matched-depths}\\
\mathcal J(\bzeta_\perp,u_\parallel,c)
&=\mathcal I\big((\bz_\perp,z_{\parallel,c}),\bzeta_\perp;c\big),
\label{eq:Jdef}
\end{align}
and
\begin{equation}
m(u_\parallel,c)^2
=\frac{u_\parallel}{U_0(y_{\parallel,c})}
=\frac{u_\parallel}
{\big(U_0\circ T_0^{-1}\circ T_c\circ U_c^{-1}\big)(u_\parallel)}.
\label{eq:m}
\end{equation}
Thus
\begin{equation}
U_c(z_{\parallel,c})=u_\parallel,\qquad
T_0(y_{\parallel,c})=T_c(z_{\parallel,c}),\qquad
U_0(y_{\parallel,c})=\frac{u_\parallel}{m^2}.
\label{eq:matching}
\end{equation}
The dimensions of $u_\parallel$ are length squared divided by time.
For a nonconstant profile, $m$ is an integrated mismatch, not a
pointwise ratio of velocities. Both matched depths depend on $c$.

\subsection{Step 4: substitute the data and the Green's functions}

Substitution of \eqref{eq:data}, \eqref{eq:G0}, and \eqref{eq:Gc} into
\eqref{eq:image} gives, before any bandwidth or depth localization,
\begin{align}
\mathcal J(\bzeta_\perp,u_\parallel,c)
\simeq{}&\int_\R\dd\omega\int_{\R^2\times\R^+}\dd\by\,
\rho(\by)\frac{1}{B}\widehat f\!\left(\frac{\omega-\omega_0}{B}\right)
\frac{c_0(0)c_0(y_\parallel)c(0)c(z_{\parallel,c})}
{256\pi^4U_0(y_\parallel)^2u_\parallel^2}
e^{2i\omega[T_0(y_\parallel)-T_c(z_{\parallel,c})]}
\nonumber\\
&\times\int_{\R^2}\dd\bx_n\,
\exp\!\left\{-\frac{|\bx_n|^2}{a^2}
+\frac{i\omega}{2}\left[
\frac{|\bx_n-\by_\perp|^2}{U_0(y_\parallel)}
-\frac{|\bx_n-\bz_\perp-\bzeta_\perp/2|^2}{u_\parallel}
\right]\right\}
\nonumber\\
&\times\int_{\R^2}\dd\bx_{n'}\,
\exp\!\left\{-\frac{|\bx_{n'}|^2}{a^2}
+\frac{i\omega}{2}\left[
\frac{|\bx_{n'}-\by_\perp|^2}{U_0(y_\parallel)}
-\frac{|\bx_{n'}-\bz_\perp+\bzeta_\perp/2|^2}{u_\parallel}
\right]\right\}.
\label{eq:J-before-localization}
\end{align}
The prefactor is the product of four Green's function amplitudes;
this explains the power $256\pi^4$ and both squared longitudinal
denominators.

\subsection{Step 5: integrate frequency under explicit conditions}

We first replace $\omega$ by $\omega_0$ only in the transverse phases.
A sufficient condition, in addition to $B\ll\omega_0$, is
\begin{equation}
B\left|
\frac{|\bx_n-\by_\perp|^2}{2U_0(y_\parallel)}
-\frac{|\bx_n-\bz_\perp\mp\bzeta_\perp/2|^2}{2u_\parallel}
\right|\ll1,
\label{eq:point-band-condition}
\end{equation}
for the contributing source and receiver positions and offsets.
Indeed, the phase error is the bracket in
\eqref{eq:point-band-condition} multiplied by $\omega-\omega_0$.
The axial phase must retain its full frequency dependence. By
\eqref{eq:fourier} and the change of variable
$\omega\mapsto(\omega-\omega_0)/B$,
\begin{align}
&\int_\R\frac{\dd\omega}{B}
\widehat f\!\left(\frac{\omega-\omega_0}{B}\right)
e^{2i\omega[T_0(y_\parallel)-T_c(z_{\parallel,c})]}
\nonumber\\
&\qquad=2\pi
f\!\left(2B[T_c(z_{\parallel,c})-T_0(y_\parallel)]\right)
e^{2i\omega_0[T_0(y_\parallel)-T_c(z_{\parallel,c})]}.
\label{eq:frequency-integral}
\end{align}
This accounts for the factor $2\pi$ and for the sign of the argument
of $f$.

\subsection{Step 6: justify freezing the transverse coefficients}

The pulse envelope in \eqref{eq:frequency-integral} selects
\begin{equation}
|T_0(y_\parallel)-T_0(y_{\parallel,c})|\lesssim B^{-1}.
\label{eq:time-window}
\end{equation}
If the profile changes little in this time window, Taylor expansion gives
\begin{equation}
T_0(y_\parallel)-T_0(y_{\parallel,c})
=\frac{y_\parallel-y_{\parallel,c}}{c_0(y_{\parallel,c})}
+O\!\left(
|y_\parallel-y_{\parallel,c}|^2
\sup\frac{|c_0'|}{c_0^2}\right),
\end{equation}
where the supremum is taken between the two depths. Hence the selected
depth range has size $c_0(y_{\parallel,c})/B$.

For all contributing depths in \eqref{eq:time-window}, we assume
\begin{gather}
\left|\frac{c_0(y_\parallel)}{c_0(y_{\parallel,c})}-1\right|\ll1,
\qquad
\left|\frac{U_0(y_\parallel)}{U_0(y_{\parallel,c})}-1\right|\ll1,
\label{eq:freeze-amplitude}\\
\omega_0|\bx_n-\by_\perp|^2
\left|\frac{1}{U_0(y_\parallel)}
-\frac{1}{U_0(y_{\parallel,c})}\right|\ll1,
\label{eq:freeze-phase}
\end{gather}
and the same phase condition for $\bx_{n'}$.
The last condition is stronger than merely requiring a small relative
change in $U_0$: it controls an oscillatory phase.
When first-order estimates are valid, sufficient local size conditions
are
\begin{equation}
\frac{|c_0'(y_{\parallel,c})|}{B}\ll1,\qquad
\frac{c_0(y_{\parallel,c})^2}{B\,U_0(y_{\parallel,c})}\ll1,\qquad
\frac{\omega_0 c_0(y_{\parallel,c})^2|\bx_n-\by_\perp|^2}
{B\,U_0(y_{\parallel,c})^2}\ll1.
\label{eq:freeze-scales}
\end{equation}
These estimates clarify why narrow relative bandwidth does not, by
itself, justify axial localization.

Under these conditions we replace $U_0(y_\parallel)$ by
$U_0(y_{\parallel,c})=u_\parallel/m^2$ in the transverse phases, and
$c_0(y_\parallel)/U_0(y_\parallel)^2$ by
$c_0(y_{\parallel,c})m^4/u_\parallel^2$ in the amplitude.
We retain $T_0(y_\parallel)$ exactly in the axial phase and envelope,
and retain the full reflectivity $\rho(\by_\perp,y_\parallel)$.

Combining these replacements with \eqref{eq:frequency-integral} yields
\begin{align}
\mathcal J(\bzeta_\perp,u_\parallel,c)
\simeq{}&
\frac{c_0(0)c_0(y_{\parallel,c})c(0)c(z_{\parallel,c})m^4}
{128\pi^3u_\parallel^4}
\int_{\R^2\times\R^+}\dd\by\,\rho(\by)
f\!\left(2B[T_c(z_{\parallel,c})-T_0(y_\parallel)]\right)
\nonumber\\
&\times e^{2i\omega_0[T_0(y_\parallel)-T_c(z_{\parallel,c})]}
\nonumber\\
&\times\int_{\R^2}\dd\bx_n\,
\exp\!\left\{-\frac{|\bx_n|^2}{a^2}
+\frac{i\omega_0}{2u_\parallel}
\left[m^2|\bx_n-\by_\perp|^2
-|\bx_n-\bz_\perp-\bzeta_\perp/2|^2\right]\right\}
\nonumber\\
&\times\int_{\R^2}\dd\bx_{n'}\,
\exp\!\left\{-\frac{|\bx_{n'}|^2}{a^2}
+\frac{i\omega_0}{2u_\parallel}
\left[m^2|\bx_{n'}-\by_\perp|^2
-|\bx_{n'}-\bz_\perp+\bzeta_\perp/2|^2\right]\right\}.
\label{eq:J-frozen}
\end{align}

\subsection{Step 7: evaluate both Gaussian integrals}

Set, with the original notation,
\begin{equation}
\alpha(u_\parallel,c)=\frac{\omega_0a^2}{2u_\parallel}(m^2-1).
\label{eq:alpha}
\end{equation}
At fixed $(u_\parallel,c)$, write $\alpha=\alpha(u_\parallel,c)$ and
$m=m(u_\parallel,c)$ in the calculations. Expanding the transverse
quadratics gives
\begin{align}
&m^2|\bx-\by_\perp|^2
-|\bx-\bz_\perp\mp\bzeta_\perp/2|^2
\nonumber\\
&\qquad=(m^2-1)|\bx|^2
-2\bx\cdot\big(m^2\by_\perp-\bz_\perp\mp\bzeta_\perp/2\big)
+m^2|\by_\perp|^2-|\bz_\perp\pm\bzeta_\perp/2|^2.
\label{eq:expand-quadratic}
\end{align}
The coefficient of $-|\bx|^2$ is $(1-i\alpha)/a^2$, whose real
part is positive. Completing the square in each transverse coordinate
gives the two-dimensional identity, for real $\bk$,
\begin{equation}
\int_{\R^2}
\exp\!\left[-\frac{1-i\alpha}{a^2}|\bx|^2
-\frac{i\omega_0}{u_\parallel}\bx\cdot\bk\right]\dd\bx
=\frac{\pi a^2}{1-i\alpha}
\exp\!\left[-\frac{\omega_0^2a^2|\bk|^2}
{4u_\parallel^2(1-i\alpha)}\right].
\label{eq:gaussian-point}
\end{equation}
For clarity, the completion in the exponent reads
\begin{align}
&-\frac{1-i\alpha}{a^2}|\bx|^2
-\frac{i\omega_0}{u_\parallel}\bx\cdot\bk
\nonumber\\
&\quad=-\frac{1-i\alpha}{a^2}
\sum_{j=1}^2\left(x_j+
\frac{i\omega_0a^2 k_j}{2u_\parallel(1-i\alpha)}\right)^2
-\frac{\omega_0^2a^2|\bk|^2}{4u_\parallel^2(1-i\alpha)}.
\end{align}
The shifted Gaussian integrals are evaluated by contour translation;
their decay is ensured by the positive real part above.

Apply \eqref{eq:gaussian-point} with
$\bk=m^2\by_\perp-\bz_\perp\mp\bzeta_\perp/2$ and restore the
constant terms in \eqref{eq:expand-quadratic}. The product of the two
array integrals is
\begin{align}
&\frac{\pi^2a^4}{(1-i\alpha)^2}
\exp\!\left[-\frac{\omega_0^2a^2}
{4u_\parallel^2(1-i\alpha)}
\left(2|m^2\by_\perp-\bz_\perp|^2
+\frac{|\bzeta_\perp|^2}{2}\right)\right]
\nonumber\\
&\qquad\times
\exp\!\left[\frac{i\omega_0}{2u_\parallel}
\left(2m^2|\by_\perp|^2-2|\bz_\perp|^2
-\frac{|\bzeta_\perp|^2}{2}\right)\right].
\label{eq:gaussian-product}
\end{align}
Here we used explicitly
\begin{align}
|m^2\by_\perp-\bz_\perp-\bzeta_\perp/2|^2
+|m^2\by_\perp-\bz_\perp+\bzeta_\perp/2|^2
&=2|m^2\by_\perp-\bz_\perp|^2+\frac{|\bzeta_\perp|^2}{2},\\
|\bz_\perp+\bzeta_\perp/2|^2
+|\bz_\perp-\bzeta_\perp/2|^2
&=2|\bz_\perp|^2+\frac{|\bzeta_\perp|^2}{2}.
\end{align}

\subsection{Step 8: isolate the offset dependence}

All terms containing $\bzeta_\perp$ in
\eqref{eq:gaussian-product} are independent of $\by$. Thus
\begin{equation}
\mathcal J(\bzeta_\perp,u_\parallel,c)
\simeq\mathcal J(\mathbf 0,u_\parallel,c)
\Psi(\bzeta_\perp,u_\parallel,c),
\label{eq:factor-point}
\end{equation}
where the zero-offset factor is
\begin{align}
\mathcal J(\mathbf0,u_\parallel,c)
\simeq{}&\frac{a^4c_0(0)c_0(y_{\parallel,c})c(0)c(z_{\parallel,c})m^4}
{128\pi u_\parallel^4[1-i\alpha(u_\parallel,c)]^2}
\int_{\R^2\times\R^+}\dd\by\,\rho(\by)
\nonumber\\
&\times f\!\left(2B[T_c(z_{\parallel,c})-T_0(y_\parallel)]\right)
e^{2i\omega_0[T_0(y_\parallel)-T_c(z_{\parallel,c})]}
\nonumber\\
&\times\exp\!\left[
-\frac{\omega_0^2a^2|m^2\by_\perp-\bz_\perp|^2}
{2u_\parallel^2[1-i\alpha(u_\parallel,c)]}
+\frac{i\omega_0}{u_\parallel}
\big(m^2|\by_\perp|^2-|\bz_\perp|^2\big)\right],
\label{eq:J0-point}
\end{align}
and the offset factor is
\begin{align}
\Psi(\bzeta_\perp,u_\parallel,c)
&=\exp\!\left[
-\frac{\omega_0^2a^2|\bzeta_\perp|^2}
{8u_\parallel^2[1-i\alpha(u_\parallel,c)]}
-\frac{i\omega_0|\bzeta_\perp|^2}{4u_\parallel}\right]
\label{eq:Psi-complex}\\
&=\exp\!\left[-\frac{|\bzeta_\perp|^2}{2\ell(u_\parallel,c)^2}
-i\gamma(u_\parallel,c)|\bzeta_\perp|^2\right].
\label{eq:Psi-point}
\end{align}
Using $(1-i\alpha)^{-1}=(1+i\alpha)/(1+\alpha^2)$, we find
\begin{align}
\ell(u_\parallel,c)^2
&=\frac{4u_\parallel^2}{\omega_0^2a^2}
\big[1+\alpha(u_\parallel,c)^2\big]
\nonumber\\
&=\frac{4u_\parallel^2}{\omega_0^2a^2}
+a^2\big(m(u_\parallel,c)^2-1\big)^2,
\label{eq:ell-point}\\
\gamma(u_\parallel,c)
&=\frac{\omega_0}{4u_\parallel}
+\frac{\omega_0^2a^2\alpha(u_\parallel,c)}
{8u_\parallel^2[1+\alpha(u_\parallel,c)^2]}
\nonumber\\
&=\frac{\omega_0}{4u_\parallel}
\frac{1+m^2(m^2-1)\omega_0^2a^4/(4u_\parallel^2)}
{1+(m^2-1)^2\omega_0^2a^4/(4u_\parallel^2)}.
\label{eq:gamma-point}
\end{align}
These are the original width and phase expressions, now accompanied by
the hypotheses needed to factor them out of the depth integral.
Without the freezing step, the offset dependence remains coupled to
$y_\parallel$, so independence from $\rho$ does not follow.

\subsection{Step 9: check the homogeneous limit}

For constant velocities,
\begin{equation}
T_c(z_\parallel)=\frac{z_\parallel}{c},\quad
U_c(z_\parallel)=cz_\parallel,\quad
z_{\parallel,c}=\frac{u_\parallel}{c},\quad
y_{\parallel,c}=\frac{u_\parallel c_0}{c^2},\quad
m=\frac{c}{c_0}.
\end{equation}
The amplitude factor in \eqref{eq:J0-point} reduces to
\begin{equation}
c_0(0)c_0(y_{\parallel,c})c(0)c(z_{\parallel,c})m^4
=c_0^2c^2m^4=c_0^4m^6.
\end{equation}
Also,
\begin{equation}
|m^2\by_\perp-\bz_\perp|^2
=m^2\left|m\by_\perp-\frac{\bz_\perp}{m}\right|^2.
\end{equation}
Hence both the prefactor and the transverse exponent agree with the
homogeneous formula. No extra depth Jacobian appears, since the
integration variable $y_\parallel$ was not changed.

\section{Plane-wave data for the same smooth profile}

\subsection{Step 10: derive the spectral Green's function}

Retain the plane-wave data and Green's function definitions
\begin{align}
\widehat v(\omega,\bkappa_n,\bkappa_{n'})
&=\iint_{\R^2\times\R^2}
e^{-i\bkappa_n\cdot\bx_n-i\bkappa_{n'}\cdot\bx_{n'}}
\widehat u(\omega,\bx_n,\bx_{n'})\dd\bx_n\dd\bx_{n'},
\label{eq:plane-data}\\
\widehat H(\omega,\bkappa,\bz)
&=\int_{\R^2}
\widehat G(\omega,(\bx_\perp,0),\bz)
e^{-i\bkappa\cdot\bx_\perp}\dd\bx_\perp.
\end{align}
These definitions use unnormalized transverse Fourier transforms.
Translation invariance gives
\begin{equation}
\widehat G(\omega,(\bx_\perp,0),\by)
=\frac{1}{(2\pi)^2}\int_{\R^2}
g(\omega,\bkappa,y_\parallel)
e^{i\bkappa\cdot(\bx_\perp-\by_\perp)}\dd\bkappa,
\end{equation}
where
\begin{equation}
\partial_{y_\parallel}^2g+
\left(\frac{\omega^2}{c(y_\parallel)^2}-|\bkappa|^2\right)g
=-\delta(y_\parallel),
\qquad
\widehat H(\omega,\bkappa,\bz)
=g(\omega,\bkappa,z_\parallel)e^{-i\bkappa\cdot\bz_\perp}.
\label{eq:g-ode}
\end{equation}
For positive frequencies and propagating, nonturning components, define
the existing longitudinal wavenumber
\begin{equation}
q(\omega,\bkappa,s)
=\sqrt{\frac{\omega^2}{c(s)^2}-|\bkappa|^2}>0.
\end{equation}
Insert $g=Ae^{i\phi}$ into \eqref{eq:g-ode} away from the source:
\begin{equation}
\partial_s^2A+2i(\partial_s\phi)(\partial_sA)
+i(\partial_s^2\phi)A+
\big[q^2-(\partial_s\phi)^2\big]A=0.
\end{equation}
The outgoing choice is $\partial_s\phi=q$. Dropping $\partial_s^2A$
at leading transport order gives
\begin{equation}
2q\,\partial_s A+(\partial_s q)A=0,
\qquad \partial_s(qA^2)=0.
\end{equation}
This WKB approximation requires, in particular,
\begin{equation}
\frac{|\partial_s q|}{q^2}\ll1,
\qquad
\frac{|\partial_s^2q|}{q^3}\ll1
\label{eq:WKB-conditions}
\end{equation}
along the relevant paths; these must hold for both profiles.
At the source, the local outgoing solution is
$i e^{iq(\omega,\bkappa,0)|y_\parallel|}/[2q(\omega,\bkappa,0)]$:
its derivative jump is $-1$. Matching fixes the amplitude and yields
\begin{equation}
g(\omega,\bkappa,y_\parallel)
\simeq\frac{i}{2\sqrt{q(\omega,\bkappa,0)
q(\omega,\bkappa,y_\parallel)}}
\exp\!\left(i\int_0^{y_\parallel}q(\omega,\bkappa,s)\dd s\right).
\label{eq:WKB-g}
\end{equation}

For the effective angular aperture, assume
$\kappa_0\overline c/\omega_0\ll1$. Then
\begin{align}
q(\omega,\bkappa,s)
&=\frac{\omega}{c(s)}
-\frac{|\bkappa|^2c(s)}{2\omega}
-\frac{|\bkappa|^4c(s)^3}{8\omega^3}
+O\!\left(\frac{|\bkappa|^6\overline c^{\,5}}{\omega^5}\right),\\
\frac{1}{\sqrt{q(\omega,\bkappa,0)q(\omega,\bkappa,y_\parallel)}}
&=\frac{\sqrt{c(0)c(y_\parallel)}}{\omega}
\left[1+O\!\left(\frac{|\bkappa|^2\overline c^{\,2}}{\omega^2}\right)\right].
\end{align}
Discarding the quartic phase requires, in addition,
\begin{equation}
\frac{\kappa_0^4}{\omega_0^3}
\int_0^{y_\parallel}c(s)^3\dd s\ll1,
\label{eq:plane-paraxial-phase}
\end{equation}
and likewise for $c_0$. Therefore,
\begin{equation}
g(\omega,\bkappa,y_\parallel)
\simeq\frac{i\sqrt{c(0)c(y_\parallel)}}{2\omega}
\exp\!\left(i\omega T_c(y_\parallel)
-\frac{i|\bkappa|^2U_c(y_\parallel)}{2\omega}\right).
\label{eq:g-paraxial}
\end{equation}

\subsection{Step 11: compute the two spectral Gaussian integrals}

The plane-wave imaging function is
\begin{align}
\mathcal I_{\rm p}(\bz,\bzeta_\perp;c)
={}&\int_\R\dd\omega
\iint_{\R^2\times\R^2}\dd\bkappa_n\dd\bkappa_{n'}
e^{-(|\bkappa_n|^2+|\bkappa_{n'}|^2)/\kappa_0^2}
\overline{\widehat H(\omega,\bkappa_n,\bz_+)}
\widehat v(\omega,\bkappa_n,\bkappa_{n'})
\overline{\widehat H(\omega,\bkappa_{n'},\bz_-)}.
\end{align}
Fourier transforming \eqref{eq:data} directly gives
\begin{align}
\widehat v(\omega,\bkappa,\bkappa')
={}&\frac{1}{B}\widehat f\!\left(\frac{\omega-\omega_0}{B}\right)
\int_{\R^2\times\R^+}\dd\by\,\rho(\by)
g_0(\omega,\bkappa,y_\parallel)
g_0(\omega,\bkappa',y_\parallel)
e^{-i(\bkappa+\bkappa')\cdot\by_\perp}.
\end{align}
There are no remaining factors of $(2\pi)$ with the stated transform
convention. Consequently,
\begin{align}
\mathcal I_{\rm p}(\bz,\bzeta_\perp;c)
={}&\int_\R\dd\omega\int_{\R^2\times\R^+}\dd\by\,
\rho(\by)\frac{1}{B}\widehat f\!\left(\frac{\omega-\omega_0}{B}\right)
\nonumber\\
&\times\int_{\R^2}\dd\bkappa\,
e^{-|\bkappa|^2/\kappa_0^2}
e^{i\bkappa\cdot(\bz_\perp+\bzeta_\perp/2-\by_\perp)}
g_0(\omega,\bkappa,y_\parallel)
\overline{g(\omega,\bkappa,z_\parallel)}
\nonumber\\
&\times\int_{\R^2}\dd\bkappa'\,
e^{-|\bkappa'|^2/\kappa_0^2}
e^{i\bkappa'\cdot(\bz_\perp-\bzeta_\perp/2-\by_\perp)}
g_0(\omega,\bkappa',y_\parallel)
\overline{g(\omega,\bkappa',z_\parallel)}.
\label{eq:Ip-exact-spectral}
\end{align}
From \eqref{eq:g-paraxial}, the Green's function product in each
spectral integral is
\begin{align}
g_0(\omega,\bkappa,y_\parallel)
\overline{g(\omega,\bkappa,z_\parallel)}
\simeq{}&\frac{\sqrt{c_0(0)c_0(y_\parallel)c(0)c(z_\parallel)}}{4\omega^2}
e^{i\omega[T_0(y_\parallel)-T_c(z_\parallel)]}
\nonumber\\
&\times\exp\!\left[-\frac{i|\bkappa|^2}{2\omega}
\big(U_0(y_\parallel)-U_c(z_\parallel)\big)\right].
\end{align}
Using the two-dimensional Gaussian formula once more, its transverse
integral, without the preceding amplitude and axial phase, is
\begin{align}
&\int_{\R^2}\dd\bkappa\,
\exp\!\left[-\left(\frac{1}{\kappa_0^2}
+\frac{i[U_0(y_\parallel)-U_c(z_\parallel)]}{2\omega}\right)|\bkappa|^2
+i\bkappa\cdot(\bz_\perp\pm\bzeta_\perp/2-\by_\perp)\right]
\nonumber\\
&\quad=\frac{\pi\kappa_0^2}
{1+i\kappa_0^2[U_0(y_\parallel)-U_c(z_\parallel)]/(2\omega)}
\nonumber\\
&\qquad\times\exp\!\left[-\frac{\kappa_0^2
|\bz_\perp\pm\bzeta_\perp/2-\by_\perp|^2}
{4\{1+i\kappa_0^2[U_0(y_\parallel)-U_c(z_\parallel)]/(2\omega)\}}\right].
\label{eq:plane-gaussian}
\end{align}
Multiplication of the two integrals therefore yields the paraxial
expression before the bandwidth and depth approximations:
\begin{align}
\mathcal I_{\rm p}(\bz,\bzeta_\perp;c)
\simeq{}&\int_\R\dd\omega\int_{\R^2\times\R^+}\dd\by\,
\rho(\by)\frac{1}{B}\widehat f\!\left(\frac{\omega-\omega_0}{B}\right)
\nonumber\\
&\times\frac{\pi^2\kappa_0^4c_0(0)c_0(y_\parallel)c(0)c(z_\parallel)}
{16\omega^4
\{1+i\kappa_0^2[U_0(y_\parallel)-U_c(z_\parallel)]/(2\omega)\}^2}
\nonumber\\
&\times e^{2i\omega[T_0(y_\parallel)-T_c(z_\parallel)]}
\exp\!\left[-\frac{\kappa_0^2
(4|\by_\perp-\bz_\perp|^2+|\bzeta_\perp|^2)}
{8\{1+i\kappa_0^2[U_0(y_\parallel)-U_c(z_\parallel)]/(2\omega)\}}\right].
\label{eq:Ip-before-localization}
\end{align}

\subsection{Step 12: localize and factor the plane-wave image}

Set $z_\parallel=z_{\parallel,c}$ and define
\begin{equation}
\mathcal J_{\rm p}(\bzeta_\perp,u_\parallel,c)
=\mathcal I_{\rm p}\big((\bz_\perp,z_{\parallel,c}),\bzeta_\perp;c\big).
\end{equation}
The same axial frequency integral \eqref{eq:frequency-integral} applies.
In addition to $B\ll\omega_0$ and slow amplitude variation, sufficient
conditions for freezing the transverse spectral phase are
\begin{equation}
\frac{B\kappa_0^2}{\omega_0^2}
|U_0(y_\parallel)-u_\parallel|\ll1,
\qquad
\frac{\kappa_0^2}{\omega_0}
|U_0(y_\parallel)-U_0(y_{\parallel,c})|\ll1.
\label{eq:plane-freezing}
\end{equation}
The first controls changing $1/\omega$ to $1/\omega_0$; the second
controls changing the true depth in the transverse phase. These
conditions are imposed on the effective pulse window and angular
aperture. In particular, the second is locally ensured by
\begin{equation}
\frac{\kappa_0^2c_0(y_{\parallel,c})^2}{\omega_0B}\ll1.
\end{equation}
Apply these approximations to \eqref{eq:Ip-before-localization}, keeping
the axial phase and $\rho$ unchanged. The result is
\begin{equation}
\mathcal J_{\rm p}(\bzeta_\perp,u_\parallel,c)
\simeq\mathcal J_{\rm p}(\mathbf0,u_\parallel,c)
\Psi_{\rm p}(\bzeta_\perp,u_\parallel,c),
\label{eq:factor-plane}
\end{equation}
with
\begin{align}
\mathcal J_{\rm p}(\mathbf0,u_\parallel,c)
\simeq{}&
\frac{\pi^3\kappa_0^4c_0(0)c_0(y_{\parallel,c})c(0)c(z_{\parallel,c})}
{8\omega_0^4
\{1+i\kappa_0^2[U_0(y_{\parallel,c})-u_\parallel]/(2\omega_0)\}^2}
\nonumber\\
&\times\int_{\R^2\times\R^+}\dd\by\,\rho(\by)
f\!\left(2B[T_c(z_{\parallel,c})-T_0(y_\parallel)]\right)
\nonumber\\
&\times e^{2i\omega_0[T_0(y_\parallel)-T_c(z_{\parallel,c})]}
\exp\!\left[-\frac{\kappa_0^2|\by_\perp-\bz_\perp|^2}
{2\{1+i\kappa_0^2[U_0(y_{\parallel,c})-u_\parallel]/(2\omega_0)\}}\right],
\label{eq:J0-plane}\\
\Psi_{\rm p}(\bzeta_\perp,u_\parallel,c)
={}&\exp\!\left[-\frac{\kappa_0^2|\bzeta_\perp|^2}
{8\{1+i\kappa_0^2[U_0(y_{\parallel,c})-u_\parallel]/(2\omega_0)\}}\right].
\label{eq:Psi-plane-complex}
\end{align}
The frequency integral multiplies the prefactor in
\eqref{eq:Ip-before-localization} by $2\pi$, giving the factor
$\pi^3/8$ in \eqref{eq:J0-plane}.

Use $U_0(y_{\parallel,c})=u_\parallel/m^2$ and rationalize:
\begin{align}
&\frac{1}{1+i\kappa_0^2[U_0(y_{\parallel,c})-u_\parallel]/(2\omega_0)}
\nonumber\\
&\quad=
\frac{1+i\kappa_0^2u_\parallel(1-m^{-2})/(2\omega_0)}
{1+\kappa_0^4u_\parallel^2(1-m^{-2})^2/(4\omega_0^2)}.
\end{align}
Thus
\begin{equation}
\Psi_{\rm p}(\bzeta_\perp,u_\parallel,c)
=\exp\!\left[-\frac{|\bzeta_\perp|^2}{2\ell_{\rm p}(u_\parallel,c)^2}
-i\gamma_{\rm p}(u_\parallel,c)|\bzeta_\perp|^2\right],
\end{equation}
where
\begin{align}
\ell_{\rm p}(u_\parallel,c)^2
&=\frac{4}{\kappa_0^2}
\left[1+\frac{\kappa_0^4u_\parallel^2}{4\omega_0^2}
(1-m^{-2})^2\right],
\label{eq:ell-plane}\\
\gamma_{\rm p}(u_\parallel,c)
&=\frac{\kappa_0^4u_\parallel}{16\omega_0}
\frac{1-m^{-2}}
{1+\kappa_0^4u_\parallel^2(1-m^{-2})^2/(4\omega_0^2)}.
\label{eq:gamma-plane}
\end{align}
At the true profile, $m=1$, the point-source phase curvature is
$\omega_0/(4u_\parallel)$, whereas the plane-wave phase curvature is
zero. Their different values follow from the different imaging kernels.

\section{Focusing criterion and identifiability}

\subsection{Step 13: state what normalization and width minimization mean}

Whenever the zero-offset image is nonzero and the approximation error is
small relative to it, \eqref{eq:factor-point} implies
\begin{equation}
\frac{\mathcal J(\bzeta_\perp,u_\parallel,c)}
{\mathcal J(\mathbf0,u_\parallel,c)}
\simeq\Psi(\bzeta_\perp,u_\parallel,c),
\end{equation}
and the analogous statement holds for plane waves. The normalized
profile is independent of reflectivity only at this leading order.
If the denominator vanishes or is too small, this data ratio is
undefined or unstable. The reduced-model function $\Psi$ itself
remains well defined.

At a fixed $u_\parallel$, \eqref{eq:ell-point} and
\eqref{eq:ell-plane} show that either width is minimized exactly when
\begin{equation}
m(u_\parallel,c)=1
\quad\Longleftrightarrow\quad
U_0\!\left(T_0^{-1}\big(T_c(U_c^{-1}(u_\parallel))\big)\right)
=u_\parallel.
\label{eq:minimum-condition}
\end{equation}
This is one scalar matching condition at this value of $u_\parallel$;
it does not identify an arbitrary smooth profile from one depth.

\subsection{Step 14: correct the objective on the infinite depth axis}

In two transverse dimensions, polar integration gives
\begin{align}
\int_{\R^2}|\Psi(\bzeta_\perp,u_\parallel,c)|^2\dd\bzeta_\perp
&=\pi\ell(u_\parallel,c)^2,\\
\int_{\R^2}|\Psi_{\rm p}(\bzeta_\perp,u_\parallel,c)|^2\dd\bzeta_\perp
&=\pi\ell_{\rm p}(u_\parallel,c)^2.
\end{align}
The factor is $\pi\ell^2$, because squaring the modulus changes
$e^{-|\bzeta_\perp|^2/(2\ell^2)}$ into
$e^{-|\bzeta_\perp|^2/\ell^2}$.
At $c=c_0$, these integrals equal
$4\pi u_\parallel^2/(\omega_0^2a^2)$ and $4\pi/\kappa_0^2$,
respectively. Their integrals over $\R^+$ both diverge.
Consequently, the unmodified infinite-depth objectives do not define
the claimed unique-minimizer problems.

Keeping the original names for the objectives, a valid idealized
replacement is to subtract the known minimum at each depth before
integrating in $u_\parallel$:
\begin{align}
\mathcal L(c)
&=\int_0^\infty\left[
\int_{\R^2}|\Psi(\bzeta_\perp,u_\parallel,c)|^2\dd\bzeta_\perp
-\frac{4\pi u_\parallel^2}{\omega_0^2a^2}\right]\dd u_\parallel
\nonumber\\
&=\pi a^2\int_0^\infty
\big(m(u_\parallel,c)^2-1\big)^2\dd u_\parallel,
\label{eq:renormalized-point}\\
\mathcal L_{\rm p}(c)
&=\int_0^\infty\left[
\int_{\R^2}|\Psi_{\rm p}(\bzeta_\perp,u_\parallel,c)|^2\dd\bzeta_\perp
-\frac{4\pi}{\kappa_0^2}\right]\dd u_\parallel
\nonumber\\
&=\frac{\pi\kappa_0^2}{\omega_0^2}
\int_0^\infty u_\parallel^2
\big(1-m(u_\parallel,c)^{-2}\big)^2\dd u_\parallel.
\label{eq:renormalized-plane}
\end{align}
These are nonnegative extended-real functionals: they may be infinite
for some trial profiles, but they are zero at $c_0$. They are defined
as integrals of the displayed nonnegative differences, not as a
subtraction of two divergent integrals. They describe the reduced
offset model. Data-based versions replace $\Psi$ by the normalized
image only where that normalization is reliable.

On any fixed bounded depth-coordinate interval, the original objective
is finite for bounded smooth profiles and subtracting this
$c$-independent baseline does not change its minimizers. This is the
appropriate practical restriction, but it changes the scope of an
identifiability claim, as explained below.

\subsection{Step 15: prove the idealized uniqueness result}

\begin{proposition}
In the reduced offset model, let $c_0$ and the admissible profiles $c$
be smooth, uniformly positive, and bounded on the positive depth axis.
If the offset profiles are available for all $u_\parallel>0$, then
each functional \eqref{eq:renormalized-point} and
\eqref{eq:renormalized-plane} has the unique minimizer $c=c_0$.
\end{proposition}

\begin{proof}
At $c=c_0$, \eqref{eq:matched-depths} gives
$y_{\parallel,c}=z_{\parallel,c}$ and $m=1$, so both functionals are
zero. Conversely, if either functional vanishes, its nonnegative
integrand is zero almost everywhere. Continuity then gives
$m(u_\parallel,c)=1$ for every $u_\parallel>0$.

With $c$ fixed, regard the already defined
$z_{\parallel,c}$ and $y_{\parallel,c}$ as functions of
$u_\parallel$. Their defining relations and $m=1$ give
\begin{equation}
U_c(z_{\parallel,c})=u_\parallel,\qquad
U_0(y_{\parallel,c})=u_\parallel,\qquad
T_0(y_{\parallel,c})=T_c(z_{\parallel,c}).
\label{eq:uniqueness-identities}
\end{equation}
Differentiating the first two identities gives
\begin{equation}
\frac{\dd z_{\parallel,c}}{\dd u_\parallel}
=\frac{1}{c(z_{\parallel,c})},\qquad
\frac{\dd y_{\parallel,c}}{\dd u_\parallel}
=\frac{1}{c_0(y_{\parallel,c})}.
\label{eq:depth-derivatives}
\end{equation}
Differentiating the travel-time identity and then using
\eqref{eq:depth-derivatives} yields
\begin{equation}
\frac{1}{c_0(y_{\parallel,c})}
\frac{\dd y_{\parallel,c}}{\dd u_\parallel}
=\frac{1}{c(z_{\parallel,c})}
\frac{\dd z_{\parallel,c}}{\dd u_\parallel}
\quad\Longrightarrow\quad
\frac{1}{c_0(y_{\parallel,c})^2}
=\frac{1}{c(z_{\parallel,c})^2}.
\end{equation}
Positivity implies
\begin{equation}
c_0(y_{\parallel,c})=c(z_{\parallel,c}),
\qquad
\frac{\dd}{\dd u_\parallel}
\big(y_{\parallel,c}-z_{\parallel,c}\big)=0.
\end{equation}
Both depths tend to zero as $u_\parallel\downarrow0$, since
$T_c(0)=T_0(0)=U_c(0)=U_0(0)=0$. Their constant difference is
therefore zero:
\begin{equation}
y_{\parallel,c}=z_{\parallel,c}.
\end{equation}
It follows that $c(z_{\parallel,c})=c_0(z_{\parallel,c})$ for every
$u_\parallel>0$. Uniform positivity and boundedness ensure that
$U_c$ maps $[0,\infty)$ onto $[0,\infty)$, so the trial depth covers
the entire positive axis. Hence $c=c_0$ everywhere, including at zero
by continuity.
\end{proof}

\subsection{Step 16: distinguish the ideal theorem from a finite experiment}

The preceding proposition is an identifiability statement for the
reduced model, not an exact uniqueness theorem for finite-frequency
Helmholtz data. Its hypotheses use the whole positive depth-coordinate
axis. The paraxial and high-frequency approximations are applied in
practice only on an observation window away from the array plane.

If $m=1$ is known only on a connected observation interval, the same
differentiation shows only that
$y_{\parallel,c}-z_{\parallel,c}$ is constant there. It need not be
zero without a depth anchor. A known overlying velocity profile and
coincident matched depths at the beginning of that interval provide
such an anchor; then the proof identifies the profile in the covered
depth interval. A finite window alone does not establish global
uniqueness among all smooth profiles, and finitely many sampled depths
do not determine an unrestricted smooth function.

The Gaussian aperture and angular cutoffs have infinite mathematical
support. Their use above is understood on the effective contributing
region, where the paraxial hypotheses hold and the omitted tails are
negligible. Likewise, integrating the Gaussian offset model over all
$\R^2$ in the ideal objective presumes its significant width lies
inside that valid region. With finite measured offsets, one uses a
fixed observed offset domain and states the corresponding restricted
data assumptions.

For several transverse midpoints, one may additionally average or
integrate the reliable normalized profiles over the existing domain
$\mathcal A_\perp$, provided the depth-only background model holds
throughout the illuminated region. This improves data averaging but
does not replace missing depth information or a missing depth anchor.

\end{document}
